
\documentclass[11pt]{article}
\usepackage[a4paper,margin=2.5cm]{geometry}
\usepackage{amsmath,amssymb,bm}
\usepackage{booktabs}
\usepackage{xcolor}
\usepackage{enumitem}
\usepackage{hyperref}

\title{Spatial Clustering of Vulnerability Regimes in TwoTimescaleResilience\\
\large Continuous, Non-Threshold Feature Construction}
\author{Draft technical note}
\date{\today}

\begin{document}
\maketitle

\section{Purpose}
This note sketches a threshold-free strategy for spatial clustering of vulnerability regimes in the
TwoTimescaleResilience framework. The goal is not to cluster raw concentration fields directly, nor to
classify cells by arbitrary exceedance thresholds such as $F>1.10$. Instead, grid cells are clustered by
continuous, biologically interpretable features derived from the model response chain: compound pressure,
DEB-axis impairment, species- or archetype-weighted impairment, adaptive-margin depletion, restoring-force
reduction, amplification, seasonal timing, axis composition, and mixture-model sensitivity.

The operational chain is
\begin{equation}
C_{j,t}(g) \;\longrightarrow\; B_{j,t}(g) \;\longrightarrow\; x_{j,t}(g)
\;\longrightarrow\; E_{a,t}^{(m)}(g)
\;\longrightarrow\; Q_{q,t}^{(m)}(g)
\;\longrightarrow\; A_{q,t}^{(m)}(g)
\;\longrightarrow\; \lambda_{q,t}^{(m)}(g)
\;\longrightarrow\; F_{q,t}^{(m)}(g),
\end{equation}
where $g$ indexes grid cells, $j$ compounds, $a$ DEB-informed physiological axes,
$q$ species or species archetypes, $t$ months, and $m$ mixture-effect models.

\section{Concentration Matrices}
Let a spatial grid contain $G=n_x n_y$ cells. For each compound $j\in\{1,\ldots,J\}$ and month
$t\in\{1,\ldots,T\}$, define a concentration matrix
\begin{equation}
\mathbf{C}_{j,t}
= \begin{bmatrix}
C_{j,t}(1,1) & C_{j,t}(1,2) & \cdots & C_{j,t}(1,n_x) \\
C_{j,t}(2,1) & C_{j,t}(2,2) & \cdots & C_{j,t}(2,n_x) \\
\vdots & \vdots & \ddots & \vdots \\
C_{j,t}(n_y,1) & C_{j,t}(n_y,2) & \cdots & C_{j,t}(n_y,n_x)
\end{bmatrix}
\in \mathbb{R}_{\ge 0}^{n_y\times n_x}.
\end{equation}
Equivalently, after vectorising the grid,
\begin{equation}
\bm{c}_{j,t}=\operatorname{vec}(\mathbf{C}_{j,t})\in\mathbb{R}_{\ge0}^{G}.
\end{equation}
Stacking compounds for one month gives
\begin{equation}
\mathbf{C}_{t}
=\begin{bmatrix}
\bm{c}_{1,t}^{\top}\\
\bm{c}_{2,t}^{\top}\\
\vdots\\
\bm{c}_{J,t}^{\top}
\end{bmatrix}
\in\mathbb{R}_{\ge0}^{J\times G}.
\end{equation}
Column $g$ is the multi-compound concentration vector in cell $g$:
\begin{equation}
\bm{c}_{t}(g)=\left(C_{1,t}(g),\ldots,C_{J,t}(g)\right)^{\top}.
\end{equation}

\section{Chemical Memory and Active Stress}
For persistent or bioaccumulative compounds, ambient concentration is converted to retained internal
burden by
\begin{equation}
B_{j,t}(g)=\rho_j B_{j,t-1}(g)+(1-\rho_j)K_j C_{j,t}(g),
\end{equation}
where $\rho_j\in[0,1)$ is monthly retention and $K_j\ge0$ is an internal magnification factor.
In vectorised form,
\begin{equation}
\bm{b}_{j,t}=\rho_j\bm{b}_{j,t-1}+(1-\rho_j)K_j\bm{c}_{j,t}.
\end{equation}

Using ECOTOX-derived no-effect and effect concentrations, active stress is
\begin{equation}
x_{j,t}(g)
=\max\left(0,\frac{B_{j,t}(g)-\mathrm{NOEC}_j}{\mathrm{EC50}_j-\mathrm{NOEC}_j}\right),
\qquad \mathrm{EC50}_j>\mathrm{NOEC}_j.
\end{equation}
Let $r(j)\in\mathcal{A}$ route compound $j$ to one DEB-informed axis,
\begin{equation}
\mathcal{A}=\{\mathrm{assimilation},\mathrm{maintenance},\mathrm{growth},\mathrm{reproduction}\}.
\end{equation}
The compound-axis burden is
\begin{equation}
x_{j,a,t}(g)=x_{j,t}(g)\,\mathbb{I}\{r(j)=a\}.
\end{equation}

\section{Axis-Level Mixture Effects}
For each grid cell, month, and axis, the mixture layer maps the set
\begin{equation}
\mathcal{X}_{a,t}(g)=\{x_{j,a,t}(g):j=1,\ldots,J\}
\end{equation}
to a bounded fractional axis impairment
\begin{equation}
E_{a,t}^{(m)}(g)\in[0,1],
\end{equation}
where $m$ denotes the mixture-effect model.

\subsection{Toxic-unit summation}
The audit pressure is
\begin{equation}
X_{a,t}(g)=\sum_{j=1}^{J}x_{j,a,t}(g),
\end{equation}
and the axis impairment is
\begin{equation}
E_{a,t}^{\mathrm{TU}}(g)=\frac{X_{a,t}(g)}{1+X_{a,t}(g)}.
\end{equation}

\subsection{Independent action}
First compute compound-level effects
\begin{equation}
E_{j,a,t}(g)=\frac{x_{j,a,t}(g)}{1+x_{j,a,t}(g)}.
\end{equation}
Then combine them as
\begin{equation}
E_{a,t}^{\mathrm{IA}}(g)=1-\prod_{j:r(j)=a}\left(1-E_{j,a,t}(g)\right).
\end{equation}

\subsection{Grouped concentration addition then independent action}
Let $h(j)$ denote an empirical effect group, for example an ECOTOX effect code. For each axis $a$ and
within-axis effect group $u$, define
\begin{equation}
X_{a,u,t}(g)=\sum_{j:r(j)=a,\;h(j)=u}x_{j,a,t}(g),
\qquad
E_{a,u,t}(g)=\frac{X_{a,u,t}(g)}{1+X_{a,u,t}(g)}.
\end{equation}
Across groups,
\begin{equation}
E_{a,t}^{\mathrm{G}}(g)=1-\prod_u\left(1-E_{a,u,t}(g)\right).
\end{equation}
The grouped model is useful as a preferred default when effect-code grouping is available: same-axis,
same-group burdens are added before transformation, while distinct groups are combined as bounded
independent effects. Toxic-unit summation and independent action can then be retained as sensitivity models,
not as claims about known synergistic or antagonistic interactions.

\section{Species or Archetype Response}
For species or archetype $q$, let $w_{q,a}$ be normalized axis weights,
\begin{equation}
\sum_{a\in\mathcal{A}}w_{q,a}=1,\qquad w_{q,a}\ge0.
\end{equation}
The weighted impairment is
\begin{equation}
Q_{q,t}^{(m)}(g)=\sum_{a\in\mathcal{A}}w_{q,a}E_{a,t}^{(m)}(g).
\end{equation}
Adaptive margin is
\begin{equation}
A_{q,t}^{(m)}(g)=A_{0,q}\max\left(\epsilon_A,1-Q_{q,t}^{(m)}(g)\right),
\qquad \epsilon_A=10^{-6}.
\end{equation}
It is useful to define the continuous relative margin remaining,
\begin{equation}
R^A_{q,t,m}(g)=\frac{A_{q,t}^{(m)}(g)}{A_{0,q}}
=\max\left(\epsilon_A,1-Q_{q,t}^{(m)}(g)\right),
\end{equation}
and the corresponding continuous margin depletion,
\begin{equation}
D^A_{q,t,m}(g)=1-R^A_{q,t,m}(g).
\end{equation}
Away from the floor, $D^A_{q,t,m}(g)=Q_{q,t}^{(m)}(g)$. This avoids introducing an external threshold:
margin depletion is measured on the model's own bounded scale.

The restoring force is
\begin{equation}
\lambda_q(A)=\lambda_{\min,q}
+\left(\lambda_{\max,q}-\lambda_{\min,q}\right)
\frac{A^+}{K_{A,q}+A^+},
\qquad A^+=\max(A,0).
\end{equation}
Finally, amplification is
\begin{equation}
F_{q,t}^{(m)}(g)=\frac{\lambda_q(A_{0,q})}{\lambda_q(A_{q,t}^{(m)}(g))}.
\end{equation}

\section{Why Not Threshold Fractions?}
A feature such as
\begin{equation}
\phi_{F>1.10}(g)=\frac{1}{|\mathcal{Q}||\mathcal{T}|}\sum_{q,t}
\mathbb{I}\{F_{q,t}(g)>1.10\}
\end{equation}
reintroduces a fixed exceedance threshold. Unless $1.10$ is supplied by an external decision context,
it is arbitrary. The clustering strategy below therefore avoids threshold fractions. Regimes are defined from
continuous distributions, relative spatial position, axis composition, seasonal timing, and mixture-model
sensitivity. If a regulator or management context later supplies a threshold, threshold summaries can be
added as optional reporting layers, not as the scientific basis of regime clustering.

\section{Continuous Vulnerability-Regime Features}
Clustering should use response-derived features, not raw concentrations alone and not arbitrary exceedance
indicators. For each cell $g$, define a feature vector $\bm{z}(g)$ using continuous summaries over species
or archetypes $\mathcal{Q}$, months $\mathcal{T}$, and mixture models $\mathcal{M}$.

\subsection{Distributional response features}
For the preferred grouped model, for example,
\begin{align}
\overline{F}^{\mathrm{G}}(g)
&=\frac{1}{|\mathcal{Q}||\mathcal{T}|}\sum_{q\in\mathcal{Q}}\sum_{t\in\mathcal{T}}
F_{q,t}^{\mathrm{G}}(g),\\
F_{0.50}^{\mathrm{G}}(g)
&=\operatorname{Quantile}_{0.50}\{F_{q,t}^{\mathrm{G}}(g):q\in\mathcal{Q},t\in\mathcal{T}\},\\
F_{0.90}^{\mathrm{G}}(g)
&=\operatorname{Quantile}_{0.90}\{F_{q,t}^{\mathrm{G}}(g):q\in\mathcal{Q},t\in\mathcal{T}\},\\
F_{0.95}^{\mathrm{G}}(g)
&=\operatorname{Quantile}_{0.95}\{F_{q,t}^{\mathrm{G}}(g):q\in\mathcal{Q},t\in\mathcal{T}\},\\
F_{\max}^{\mathrm{G}}(g)
&=\max_{q\in\mathcal{Q},\;t\in\mathcal{T}}F_{q,t}^{\mathrm{G}}(g).
\end{align}
Analogous summaries can be computed for $Q_{q,t}^{\mathrm{G}}(g)$ or for relative margin remaining
$R^A_{q,t,\mathrm{G}}(g)$.

\subsection{Species or archetype heterogeneity}
Cells can differ not only in their mean response, but also in how unevenly species or archetypes respond.
A continuous heterogeneity feature is
\begin{equation}
\operatorname{sd}_{q,t}\left(F^{\mathrm{G}}\right)(g)
=\operatorname{sd}\{F_{q,t}^{\mathrm{G}}(g):q\in\mathcal{Q},t\in\mathcal{T}\}.
\end{equation}
A high value indicates that some species or archetypes respond much more strongly than others, even if the
mean response is moderate.

\subsection{Relative spatial position}
To avoid fixed cut-offs while still identifying spatially extreme cells, use rank-percentile features. For a
cell-level scalar $S(g)$, define
\begin{equation}
\Pi_S(g)=\frac{1}{G}\sum_{g'=1}^{G}\mathbb{I}\{S(g')\le S(g)\}.
\end{equation}
For example, $\Pi_{F_{\max}^{\mathrm{G}}}(g)$ is the spatial percentile rank of the cell's maximum grouped-model
amplification. This says whether a cell is high relative to the study domain without declaring any absolute
amplification threshold biologically special.

\subsection{Mixture-model sensitivity}
Mixture-effect assumptions should enter as continuous deltas, not as categorical claims about true interactions.
Useful features include
\begin{align}
\Delta F_{\mathrm{IA}-\mathrm{TU}}(g)
&=\max_{q,t}\left(F_{q,t}^{\mathrm{IA}}(g)-F_{q,t}^{\mathrm{TU}}(g)\right),\\
\Delta F_{\mathrm{G}-\mathrm{TU}}(g)
&=\max_{q,t}\left(F_{q,t}^{\mathrm{G}}(g)-F_{q,t}^{\mathrm{TU}}(g)\right),\\
\delta F_{\mathrm{IA}-\mathrm{TU}}(g)
&=\max_{q,t}\frac{F_{q,t}^{\mathrm{IA}}(g)-F_{q,t}^{\mathrm{TU}}(g)}{F_{q,t}^{\mathrm{TU}}(g)},\\
\delta F_{\mathrm{G}-\mathrm{TU}}(g)
&=\max_{q,t}\frac{F_{q,t}^{\mathrm{G}}(g)-F_{q,t}^{\mathrm{TU}}(g)}{F_{q,t}^{\mathrm{TU}}(g)}.
\end{align}
These features identify where same-axis overlap and effect-code grouping materially affect inferred vulnerability.

\subsection{Axis composition and axis diversity}
Let
\begin{equation}
\overline{E}_a^{\mathrm{G}}(g)=\frac{1}{|\mathcal{T}|}\sum_{t\in\mathcal{T}}E_{a,t}^{\mathrm{G}}(g).
\end{equation}
The dominant axis is
\begin{equation}
a^*(g)=\arg\max_{a\in\mathcal{A}}\overline{E}_a^{\mathrm{G}}(g).
\end{equation}
To measure whether impairment is concentrated in one axis or distributed across several axes, define
\begin{equation}
p_a(g)=\frac{\overline{E}_a^{\mathrm{G}}(g)}{\sum_{a'\in\mathcal{A}}\overline{E}_{a'}^{\mathrm{G}}(g)},
\end{equation}
when the denominator is positive. The axis-diversity entropy is
\begin{equation}
H_E(g)=-\sum_{a\in\mathcal{A}}p_a(g)\log p_a(g).
\end{equation}
Low $H_E$ means one axis dominates. High $H_E$ means impairment is spread across axes.

\subsection{Seasonal timing}
Seasonal timing can be represented without thresholds by recording the month of maximum response:
\begin{equation}
t^*_F(g)=\arg\max_{t\in\mathcal{T}}\operatorname{Quantile}_{0.95}
\{F_{q,t}^{\mathrm{G}}(g):q\in\mathcal{Q}\}.
\end{equation}
This distinguishes persistent chronic vulnerability from seasonal or pulse-dominated vulnerability.

\subsection{Example threshold-free feature vector}
A compact feature vector could be
\begin{equation}
\bm{z}(g)=
\begin{bmatrix}
\overline{Q}^{\mathrm{G}}(g)\\
Q_{0.50}^{\mathrm{G}}(g)\\
Q_{0.90}^{\mathrm{G}}(g)\\
Q_{0.95}^{\mathrm{G}}(g)\\
\overline{F}^{\mathrm{G}}(g)\\
F_{0.50}^{\mathrm{G}}(g)\\
F_{0.90}^{\mathrm{G}}(g)\\
F_{0.95}^{\mathrm{G}}(g)\\
F_{\max}^{\mathrm{G}}(g)\\
\operatorname{sd}_{q,t}(F^{\mathrm{G}})(g)\\
\Pi_{F_{\max}^{\mathrm{G}}}(g)\\
\Delta F_{\mathrm{IA}-\mathrm{TU}}(g)\\
\Delta F_{\mathrm{G}-\mathrm{TU}}(g)\\
\delta F_{\mathrm{IA}-\mathrm{TU}}(g)\\
\delta F_{\mathrm{G}-\mathrm{TU}}(g)\\
\overline{E}_{\mathrm{assimilation}}^{\mathrm{G}}(g)\\
\overline{E}_{\mathrm{maintenance}}^{\mathrm{G}}(g)\\
\overline{E}_{\mathrm{growth}}^{\mathrm{G}}(g)\\
\overline{E}_{\mathrm{reproduction}}^{\mathrm{G}}(g)\\
H_E(g)\\
t^*_F(g)
\end{bmatrix}.
\end{equation}
The exact feature set can be reduced after inspecting collinearity. The important design principle is that
features are continuous summaries of model behaviour, not arbitrary exceedance counts.

\section{Standardisation and Clustering}
Let $\mathbf{Z}\in\mathbb{R}^{G\times P}$ be the feature matrix whose row $g$ is $\bm{z}(g)^{\top}$.
Features with undefined values, for example entropy in cells with zero total impairment, should be assigned
well-defined neutral values or excluded before clustering. Each retained feature is standardised:
\begin{equation}
\widetilde{Z}_{g,p}=\frac{Z_{g,p}-\mu_p}{\sigma_p},
\end{equation}
where $\mu_p$ and $\sigma_p$ are the mean and standard deviation of feature $p$ over all grid cells.

A simple vulnerability-regime partition can be obtained with $K$-means:
\begin{equation}
\min_{\mathcal{C}_1,\ldots,\mathcal{C}_K}
\sum_{k=1}^{K}\sum_{g\in\mathcal{C}_k}
\left\|\widetilde{\bm{z}}(g)-\bm{\mu}_k\right\|_2^2,
\end{equation}
where $\bm{\mu}_k$ is the centroid of regime $k$. Other clustering methods can be substituted, but the
scientific point remains: cells are grouped by continuous response-regime features rather than by fixed
threshold exceedance.

\section{Interpreting Regimes}
After clustering, each cell receives a regime label
\begin{equation}
R(g)\in\{1,\ldots,K\}.
\end{equation}
Regimes should be named after their centroid features, not after arbitrary thresholds. Example labels include:
\begin{description}[leftmargin=3.6cm]
\item[Low-response regime:] Low mean and upper-quantile $Q$ and $F$.
\item[Maintenance-dominated regime:] High $\overline{E}_{\mathrm{maintenance}}$ and low axis entropy.
\item[Multi-axis impairment regime:] High axis entropy and elevated upper-quantile $Q$.
\item[Pulse-dominated regime:] High upper-envelope $F$ concentrated in a particular season, indicated by $t^*_F$.
\item[Mixture-sensitive regime:] Large continuous mixture-model deltas, indicating sensitivity to TU, IA, or grouped assumptions.
\item[Species-heterogeneous regime:] High $\operatorname{sd}_{q,t}(F)$, indicating strong differences among archetypes or species.
\end{description}

\section{Recommended Visual Outputs}
The clustering workflow should produce maps and summaries such as:
\begin{enumerate}
\item A categorical map of $R(g)$, the vulnerability-regime label.
\item Centroid bar plots showing the defining continuous features of each regime.
\item Quantile maps such as $F_{0.95}^{\mathrm{G}}(g)$ and $Q_{0.95}^{\mathrm{G}}(g)$.
\item Spatial percentile maps such as $\Pi_{F_{\max}^{\mathrm{G}}}(g)$.
\item Mixture-sensitivity maps such as $\Delta F_{\mathrm{IA}-\mathrm{TU}}(g)$ and $\Delta F_{\mathrm{G}-\mathrm{TU}}(g)$.
\item Dominant-axis maps $a^*(g)$ and axis-diversity maps $H_E(g)$.
\item Seasonal timing maps $t^*_F(g)$.
\end{enumerate}

\section{Practical Note}
For many species, full species-by-cell heatmaps are not the primary visual object. A better strategy is to
cluster cells using threshold-free summary features across archetypical species or all-species summaries,
then map regimes, quantiles, spatial percentiles, dominant axes, axis diversity, seasonal timing, and
mixture-model sensitivity. Full rasters can then be produced only for representative archetypes or top-ranked
species within each regime.

\section{Summary}
Spatial clustering can help transform high-dimensional grid outputs into interpretable vulnerability regimes,
but the clustering should not reintroduce arbitrary thresholds. The recommended approach is to construct
continuous features from the model's own quantities: bounded axis impairments, weighted impairment,
relative margin remaining, restoring-force amplification, response heterogeneity, axis composition, seasonal
timing, and mixture-effect sensitivity. Threshold summaries may be added later only when a decision context
supplies an externally justified cut-off.

\end{document}
